\documentclass[../main.tex]{subfiles}

\begin{document}

\section{Probability Theory}

Ultimately, ML models are prediction engines. In other words,
it is very crucial to understand fundamental probability theory.
Therefore, let's get started with basic terminologies and concepts.

\begin{definition}{}{}
\textbf{Probability distribution} is a process of describing
probability and possible cases.
\end{definition}

For example, using functions or tables to represents probability of
getting a certain rolling sum after rolling a fair standard die two
times is using probability distribution.

Few symbols to keep in mind are $~$ is used to represent the
following of the probability distribution, and $\propto$ is used to
illustrate direct proportionality.

\begin{definition}{}{}
\textbf{Sample space}, denoted as $\Omega$, is a set of possible
outcomes. \textbf{Probability measure} is a function that maps
such outcome events to probability in $[0, 1]$. \textbf{Random
Variables} are functions that map the outcomes to certain numbers.
\end{definition}

Notice that the sum of all measures in a sample space is $1$.
Now sometimes, probabilities depend on the previous events.

\begin{definition}{}{}
If $A$ is the event that occurs after $B$, the probability of $A$
occurring after $B$ is known as the \textbf{posterior probability},
denoted as $P(A \mid B)$. Moreover, $P(A)$ is the \textbf{prior}
probability.
\end{definition}

For instance, if $A$ is the event where you see a snail outside and
$B$ is the event of raining, then the probability $P(A \mid B)$ is
the probability of you seeing a snail outside given that it is
raining. With this in mind, we can observe the following important
theorem that will come very handy when we discuss language models.

\begin{theorem}
{The Chain Rule of Probability}
{The Chain Rule of Probability}
The chain rule of probability states the following equation.
\[
P(A, B) = P(A \mid B) P(B)
\]
\end{theorem}

The equation above is read as the probability of $A$ and $B$
is equal to the probability of $A$ given $B$ multiplied by
the probability of $B$. The direct corollary of the chain rule
is Bayes' Theorem.

\begin{theorem}
{Bayes' Theorem}
{Bayes' Theorem}
Conditional probabilities can be found with Bayes' Theorem
defined with the equation below for $P(B) \neq 0$.
\[
P(A \mid B) = \frac{P(B \mid A) P(A)}{P(B)}
\]
\end{theorem}

As you may have guessed, we can scale the chain rule above
with generalization. Consider the following generalized version
of the chain rule.

\begin{theorem}
{Generalized Chain Rule}
{Generalized Chain Rule}
The chain rule of probability states the following.
\begin{align*}
    P(x_1, x_2, \ldots, x_n)
    &= P(x_1) P(x_2 \mid x_1) P(x_3 \mid x_1, x_2) \\
        &\qquad\qquad\qquad\cdots
            P(x_n \mid x_1, x_2, \ldots, x_{n-1}) \\
    &= \prod_{i = 1}^{n} P \left(
        x_i \; \middle|\; \bigcap_{k = 1}^{i - 1} x_{k}
    \right)
\end{align*}
\end{theorem}
\begin{proof}
Deriving this generalization from the chain rule with two
variables is rather simple. We can use the definition multiple
times to multiply and telescope. Here is what I mean. (Note
that mathematicians like to use the standard set notations while
ML engineers likes to use a more ``colloquial notation''. For
the purpose of these notes, let's use the notations that ML
engineers use.)

Consider the following equations by the definition of conditional
probability.
\begin{align*}
    P(x_1) &= P(x_1) \\
    P(x_2 \mid x_1) &= \frac{P(x_1, x_2)}{P(x_1)} \\
    P(x_3 \mid x_1, x_2) &= \frac{P(x_1, x_2, x_3)}{P(x_1, x_2)} \\
    &\ \vdots \\
    P(x_n \mid x_1, \ldots, x_{n-1})
        &= \frac{P(x_1, \ldots, x_n)}{P(x_1, \ldots, x_{n-1})}
\end{align*}
Multiplying the equations, the following equation is obtained.
\begin{align*}
    P(x_1) P(x_2 &\mid x_1) P(x_3 \mid x_1, x_2)
        P(x_n \mid x_1, \ldots, x_{n-1}) \\
    &= P(x_1) \cdot \frac{P(x_1, x_2)}{P(x_1)} \cdot
        \frac{P(x_1, x_2, x_3)}{P(x_1, x_2)} \cdots
        \frac{P(x_1, \ldots, x_n)}{P(x_1, \ldots, x_{n-1})}
\end{align*}
Telescoping the right hand side, we get our final equation.
\[
P(x_1, \ldots, x_n)
= P(x_1) P(x_2 \mid x_1) P(x_3 \mid x_1, x_2)
    P(x_n \mid x_1, \ldots, x_{n-1})
\]
Alternatively, we could try going backward from the final
equation from the definition, but I found this proof a little
more satisfying.
\end{proof}



\end{document}


